\documentclass[border=0pt]{standalone}
\usepackage{fontspec}
\setmainfont{Liberation Serif}
\usepackage{unicode-math}
\setmathfont{latinmodern-math.otf}
\usepackage{tkz-euclide}
\usetikzlibrary{arrows.meta,calc,intersections,decorations.markings,angles,quotes}

% tkz-euclide assumes centimetre coordinates. Fonts and strokes use print sizes.
\tikzset{
  every picture/.style={x=1cm,y=1cm,line width=.55pt,line cap=round,line join=round,font=\fontsize{11}{13}\selectfont},
  every node/.style={inner sep=2pt},
  construction/.style={densely dashed,line width=.4pt},
  dimension/.style={<->,>=Stealth,line width=.4pt},
  vector/.style={postaction=decorate,decoration={markings,mark=at position .55 with {\arrow{Stealth[length=4pt,width=3pt]}}}},
  note/.style={font=\fontsize{9}{11}\selectfont\itshape,draw=black!65,line width=.35pt,inner sep=3pt,align=center},
  point/.style={circle,fill=black,inner sep=0pt,minimum size=1.8pt},
  axis/.style={-{Stealth[length=4pt,width=3pt]},line width=.45pt}
}
\newcommand{\NotToScale}[1]{\node[note] at #1 {Diagram is\\NOT to scale};}
% Fix the output canvas so Typst does not resize typography to fit label extents.
\newcommand{\ExamBounds}[2]{\pgfresetboundingbox\path[use as bounding box] (0,0) rectangle (#1mm,#2mm);}

\begin{document}
\begin{tikzpicture}
  % Coordinates in cm; 1 cm in the question = 0.6 cm here. C lies on the ray
  % from B at 35 degrees above BA, at distance 14cos35 -/+ sqrt(10^2 - (14sin35)^2) from B.
  \def\k{.6}
  \coordinate (A) at (3.6,1.4);
  \coordinate (B) at ($(A)+(14*\k,0)$);
  \pgfmathsetmacro{\dl}{14*cos(35)-sqrt(100-(14*sin(35))^2)}
  \pgfmathsetmacro{\dr}{14*cos(35)+sqrt(100-(14*sin(35))^2)}
  \coordinate (C1) at ($(B)+(145:\dl*\k)$);
  \coordinate (C2) at ($(B)+(145:\dr*\k)$);
  % Circle centre A radius 10 cm, shown as a faint construction arc.
  \draw[construction,black!45] ($(A)+(12:10*\k)$) arc[start angle=12,end angle=100,radius=10*\k];
  \draw (A) -- (B) -- (C2) -- cycle;
  \draw (A) -- (C1);
  \pic[draw,line width=.45pt,angle radius=11mm,"$35^\circ$",angle eccentricity=1.38] {angle=C1--B--A};
  \foreach \p in {A,B,C1,C2} \node[point] at (\p) {};
  \node[below left] at (A) {$A$};
  \node[below right] at (B) {$B$};
  \node[anchor=south west,inner sep=1.5pt] at (C1) {$C_1$};
  \node[above right] at (C2) {$C_2$};
  \node[below=1.5mm] at ($(A)!.5!(B)$) {14 cm};
  \path (A) -- node[sloped,below,pos=.58] {10 cm} (C1);
  \path (A) -- node[left=1mm,pos=.5] {10 cm} (C2);
  \NotToScale{(13.8,7.2)}
  \ExamBounds{156}{85.8}
\end{tikzpicture}
\end{document}
